% ============================================================
%  The LaTeXEditor Guide
%
%  This document is written about the editor, in the editor.
%  Open the enclosing folder with "Open project" and press
%  Typeset (Cmd+Enter) to build it.
%
%  Note: the editor runs a SINGLE compiler pass, so press
%  Typeset twice the first time — the table of contents and
%  cross-references need the .aux file from the first run.
%
%  Packages used are deliberately limited to ones present in
%  essentially every LaTeX installation, so this compiles on a
%  minimal TeX Live/MiKTeX as well as a full one.
% ============================================================
\documentclass[11pt,a4paper]{article}

\usepackage[utf8]{inputenc}
\usepackage[T1]{fontenc}
\usepackage[margin=2.6cm]{geometry}
\usepackage{amsmath,amssymb}
\usepackage{booktabs}
\usepackage{enumitem}
\usepackage{xcolor}
\usepackage{listings}
\usepackage[hidelinks]{hyperref}

% parskip's job, done with core LaTeX so no package is needed.
\setlength{\parindent}{0pt}
\setlength{\parskip}{0.65em}

% ---- Colours picked to match the editor's own palette -------
\definecolor{accent}{HTML}{A85A2C}
\definecolor{codebg}{HTML}{F4F0E6}
\definecolor{codekw}{HTML}{A85A2C}
\definecolor{codecmt}{HTML}{8D8373}

% ---- A listing style for LaTeX snippets ---------------------
% No escapechar: everything inside a listing is taken literally,
% which is what you want when the samples are themselves LaTeX.
\lstdefinestyle{tex}{
  language=[LaTeX]TeX,
  basicstyle=\ttfamily\small,
  keywordstyle=\color{codekw}\bfseries,
  commentstyle=\color{codecmt}\itshape,
  stringstyle=\color{black},
  backgroundcolor=\color{codebg},
  frame=single,
  rulecolor=\color{codecmt!40},
  framesep=6pt,
  xleftmargin=6pt,
  xrightmargin=6pt,
  breaklines=true,
  showstringspaces=false,
  columns=fullflexible,
  keepspaces=true,
}
\lstset{style=tex}

% A shorthand for naming UI elements consistently.
\newcommand{\ui}[1]{\textsf{\textbf{#1}}}
\newcommand{\kbd}[1]{\texttt{#1}}

\title{\textbf{LaTeXEditor}\\[0.3em]
       \large A guide to writing, typesetting and revising with it}
\author{Project documentation}
\date{\today}

\begin{document}

\maketitle

\begin{abstract}
\noindent
LaTeXEditor is a desktop LaTeX environment: a bundled Monaco editor with a
LaTeX grammar, a local typesetting toolchain, a PDF preview, Git integration,
and an AI assistant that edits your document rather than merely describing
changes. This guide explains what each part does and works through the tasks
you will actually perform --- opening a project, typesetting, reading errors,
and using AI on a selection. It is itself a LaTeXEditor project, so you can
open this folder and typeset it as your first exercise.
\end{abstract}

\tableofcontents

\section{What the editor is}
\label{sec:overview}

LaTeXEditor is a desktop application. It does not upload your documents
anywhere, and it does not typeset in the cloud: it drives the \LaTeX{}
distribution already installed on your machine, and everything it needs to
render --- the code editor, the PDF viewer --- is bundled into the app rather
than fetched from a CDN. Open it on a plane and it behaves exactly as it does
at your desk.

\subsection{The window}

The layout is four regions, which this diagram shows in the positions they
actually occupy:

\begin{center}
\setlength{\fboxsep}{6pt}
\begin{tabular}{c}
  \fbox{\parbox{0.82\textwidth}{\centering \textbf{Toolbar} \\ \small
        document name \quad$\cdot$\quad compiler \quad$\cdot$\quad Typeset}} \\[4pt]
  \begin{tabular}{@{}c@{\hspace{4pt}}c@{\hspace{4pt}}c@{}}
    \fbox{\parbox[t][2.4cm]{0.22\textwidth}{\centering \textbf{Sidebar} \\[3pt] \small
          Files \\ Components \\ AI \\ Git \\ Outline}} &
    \fbox{\parbox[t][2.4cm]{0.30\textwidth}{\centering \textbf{Editor} \\[3pt] \small
          your \texttt{.tex} source}} &
    \fbox{\parbox[t][2.4cm]{0.24\textwidth}{\centering \textbf{Preview} \\[3pt] \small
          the typeset PDF}} \\
  \end{tabular} \\[4pt]
  \fbox{\parbox{0.82\textwidth}{\centering \textbf{Issues} \\ \small
        errors and warnings from the last run}} \\
\end{tabular}
\end{center}

Every divider between those regions can be dragged, and the sidebar, preview
and issues panel can each be hidden --- from the command palette, or by
clicking the \ui{$\times$} on the issues header. The editor is given a minimum
width that the preview is not allowed to squeeze past, so no amount of dragging
or window-resizing can collapse the thing you are writing in.

\subsection{What makes it different}

Three choices are worth knowing about up front, because they shape how the rest
of the tool behaves.

\begin{description}[style=nextline, leftmargin=0pt, labelindent=0pt, labelwidth=0pt, font=\normalfont\itshape]

  \item[The AI edits the document.]
  Most editors bolt on a chat window that produces text you then copy by hand.
  Here you press \kbd{Cmd+K} in the paragraph you want changed, say what you
  want, and get a diff that tints the words that actually differ. Accept and it
  is written back into exactly the range it came from.
  Section~\ref{sec:ai} is devoted to this.

  \item[It knows your project, not just your file.]
  When the assistant generates a table, it is told which packages your preamble
  loads; when it writes a citation, it is given the keys that actually exist in
  your \texttt{.bib}; when it fixes an error, it is handed the error. You are
  never asked to paste a log.

  \item[It typesets locally, once.]
  A single pass of \texttt{pdflatex} (or \texttt{xelatex}, or
  \texttt{lualatex}) per press. This is fast and predictable, and it has one
  consequence you must know about --- see Section~\ref{sec:passes}.

\end{description}

\subsection{Requirements}

You need a \LaTeX{} distribution: \textbf{TeX Live} on Linux and macOS,
\textbf{MacTeX} on macOS, or \textbf{MiKTeX} on Windows. On launch the editor
looks for \texttt{pdflatex}, \texttt{xelatex}, \texttt{lualatex} and
\texttt{latexmk}, and lists whichever it finds in the toolbar's compiler menu.

It does not rely on your \texttt{PATH} alone, which matters more than it
sounds: an app launched from the Dock or Finder inherits a minimal environment
--- roughly \path{/usr/bin:/bin:/usr/sbin:/sbin} --- so a TeX installation that
works perfectly in a terminal can be invisible to a GUI app. The editor
therefore also checks the standard install locations directly
(\path{/Library/TeX/texbin}, Homebrew, MacPorts, per-user \path{~/.local/bin},
and versioned TeX Live trees) and passes the compiler's own directory down to
it, since \texttt{latexmk} shells out to \texttt{pdflatex} in turn.

If no compiler is found, the editor says so in a banner with a
\ui{Re-check} button --- install a distribution and press it, no restart
needed. Editing works regardless; only typesetting is blocked.

\section{Getting started}
\label{sec:start}

\subsection{Your first typeset}

Do this now, with the document you are reading.

\begin{enumerate}[leftmargin=*]
  \item Click \ui{Open project} in the toolbar and choose the folder
        containing this file (\path{demo-project/le-guide}).
  \item In the \ui{Files} panel, click \texttt{main.tex}.
  \item Press \kbd{Cmd+Enter}, or click \ui{Typeset}.
\end{enumerate}

The PDF appears on the right. The \ui{Issues} panel along the bottom fills with
anything the run reported.

\subsection{You will need to press it twice}
\label{sec:passes}

The first build of this guide produces a document whose table of contents is
empty and whose cross-references read \texttt{??}. Press \ui{Typeset} again and
both fill in.

That is not a bug, and it is not specific to this editor --- it is how \LaTeX{}
works. Cross-references, the table of contents and page numbers are written to
an auxiliary file on one pass and read back on the next. Tools like
\texttt{latexmk} hide this by re-running the compiler until the output settles.
LaTeXEditor runs \emph{one} pass per press, which is faster and makes the error
output easier to read, but leaves the second press to you.

Two practical consequences:

\begin{itemize}[leftmargin=*]
  \item \textbf{Cross-references and the ToC need two presses} after they
        change. Just press \ui{Typeset} again.
  \item \textbf{BibTeX and Biber are not run.} A \verb|\bibliography{refs}|
        will never resolve, because the editor does not invoke a second
        program. This guide therefore uses a \texttt{thebibliography}
        environment written directly into \texttt{main.tex}, which resolves in
        a single pass. If your project needs full BibTeX, run it once in a
        terminal, then carry on typesetting here.
\end{itemize}

If you would rather have the loop handled for you, select \texttt{latexmk} in
the compiler menu --- it performs its own re-runs internally.

\subsection{Where the output goes}

Everything the compiler emits --- \texttt{.pdf}, \texttt{.aux}, \texttt{.log},
\texttt{.toc} --- is written to a \texttt{build/} folder beside your project,
never next to your sources. Your project folder stays readable.

\subsection{Which file gets typeset}

Not, in general, the one you are looking at. Before each run the editor scans
the project for the \texttt{.tex} file containing \verb|\documentclass| and
typesets \emph{that}.

This is what you want in a multi-file project. Editing
\texttt{sections/03-writing.tex} and pressing \ui{Typeset} builds
\texttt{main.tex}, which pulls the section back in via \verb|\input|. You never
have to navigate back to the root file to build.

\subsection{Saving}

There is no save button and you rarely need \kbd{Cmd+S}. The editor writes your
file one second after you stop typing. A dot beside the filename in the toolbar
means there are changes not yet on disk; it disappears when they are. Switching
files flushes any pending write immediately, so edits made in the last second
before you click away are not lost.

Turn on \ui{Auto} in the toolbar and every save also triggers a typeset --- a
tight write-and-see loop, at the cost of a compiler run every time you pause.

\subsection{Starting from a template}

\ui{New from template} offers a gallery: Basic Article, IEEE Conference Paper,
ACM Paper, Thesis/Dissertation, Beamer Presentation, CV/Resume, Cover Letter and
Homework/Problem Set. Choosing one writes a complete, compiling project into
your current folder and opens its main file. It is a faster start
than an empty buffer, and a reasonable way to see how a document class you have
not used before is normally set up.

\section{Writing}
\label{sec:writing}

\subsection{The editor}

The editing surface is Monaco --- the engine behind VS Code --- with a LaTeX
grammar written for this app. Commands, env names, maths and comments
are all coloured differently, \verb|\begin|/\verb|\end| blocks fold, and
\kbd{Cmd+/} comments the current line.

The editor's colours come from the same palette as the rest of the interface,
so switching between the dark and light themes in \ui{Settings} restyles the
editor along with everything else, rather than leaving a dark rectangle in a
light window.

\subsection{Snippets}

Start typing a construct's name and the completion list appears; press
\kbd{Tab} to move between the placeholders it leaves behind. Typing
\texttt{figure} and accepting the suggestion gives you:

\begin{lstlisting}
\begin{figure}[htbp]
	\centering
	\includegraphics[width=0.8\textwidth]{path}
	\caption{caption}
	\label{fig:label}
\end{figure}
\end{lstlisting}

with the cursor on \texttt{0.8\textbackslash textwidth} and \kbd{Tab} stepping
through the rest. The full set covers \texttt{begin}, \texttt{frac},
\texttt{sqrt}, \texttt{sum}, \texttt{int}, \texttt{lim}, \texttt{prod},
\texttt{section}, \texttt{subsection}, \texttt{itemize}, \texttt{enumerate},
\texttt{figure}, \texttt{table}, \texttt{align}, \texttt{matrix},
\texttt{cases}, \texttt{cite}, \texttt{ref}, and the text-formatting commands.

\subsection{The component library}

Snippets are skeletons; components are filled-in blocks. Open \ui{Components}
in the sidebar, pick one, complete the form, and press \ui{Insert} to place the
generated \LaTeX{} at your cursor.

\begin{table}[htbp]
  \centering
  \caption{What the component library offers}
  \label{tab:components}
  \begin{tabular}{ll}
    \toprule
    \textbf{Category} & \textbf{Examples} \\
    \midrule
    Mathematics        & numbered and aligned equations, matrices, cases, theorems \\
    Tables             & basic and publication tables with \texttt{booktabs} rules \\
    Figures            & single figures, side-by-side subfigures, Ti\emph{k}Z diagrams \\
    Lists              & itemised and enumerated lists \\
    Structure          & sections, abstracts \\
    Code               & source listings, algorithm and pseudocode blocks \\
    Bibliography       & article and book Bib\TeX{} entries \\
    \bottomrule
  \end{tabular}
\end{table}

Use the library when you want a correct starting point for something fiddly ---
a \texttt{booktabs} table with the rules in the right places, or a subfigure
pair with aligned captions --- and snippets when you already know the shape and
just want the keystrokes saved.

\subsection{Moving around a document}

The \ui{Outline} panel lists every heading in the open file and jumps to it on
click. It reads \verb|\part| through \verb|\subparagraph|, understands starred
forms like \verb|\section*|, and ignores anything commented out --- so a
heading you have temporarily disabled with a \texttt{\%} does not clutter it.

For longer jumps, the command palette (\kbd{Cmd+Shift+P}) exposes go-to-line,
find, find-and-replace, comment toggling, line duplication and movement, and
the editor zoom controls, and the full shortcut list --- search it for
\emph{keyboard shortcuts}. \kbd{Cmd+K} is not the way in: that key opens the
inline assistant, Section~\ref{sec:inline}.

\subsection{Reading errors}

When a run fails, the \ui{Issues} panel lists what went wrong: errors first,
then warnings, then overfull and underfull boxes. Each row gives the message,
the offending source line where the log provides it, and the file and line
number.

Click a row and the editor opens that file and puts your cursor on that line.

The distinction is worth respecting. An \emph{error} stopped the compiler and
must be fixed. A \emph{warning} --- an undefined reference, a missing citation
--- produced a PDF anyway, but usually one with \texttt{??} in it. A
\emph{badbox} is purely typographic: a line \LaTeX{} could not break neatly.
Badboxes are not clickable, because there is generally nothing at that line to
fix; they are a signal to reword a sentence or let it go.

Remember Section~\ref{sec:passes}: a fresh project reports undefined references
on its first run and resolves them on the second. Do not go hunting for a bug
that a second press of \ui{Typeset} will fix.

\section{The AI assistant}
\label{sec:ai}

The assistant is built around one idea: it should change your document, not
describe a change and leave you to transcribe it. Everything below follows from
that.

\subsection{Setting it up}

Open \ui{Settings}, choose a provider, paste its key, and press
\ui{Test Connection}. Supported providers are OpenAI, Anthropic, DeepSeek,
OpenRouter, Grok, OpenCode Go, and Ollama --- the last running locally and
needing no key at all, which is the option to reach for if you would rather
your draft never left the machine.

Keys are stored in \texttt{\textasciitilde/.latex-editor/settings.json}.

\subsection{Model and speed}
\label{sec:ai-model}

Skills work on short passages, so what matters is latency, not raw capability.
The thing that dominates latency is whether the model \emph{deliberates} before
answering --- and most current models do, at length.

Proofreading one section of this guide, identical text and prompt each time:

\begin{table}[htbp]
  \centering
  \caption{Measured on this document, via OpenCode Go}
  \label{tab:models}
  \begin{tabular}{lrr}
    \toprule
    \textbf{Configuration} & \textbf{Time} & \textbf{Thinking} \\
    \midrule
    \texttt{deepseek-v4-pro}                        & 150\,s & 40{,}700 chars \\
    \texttt{deepseek-v4-flash}                      & 45\,s  & 22{,}600 chars \\
    \texttt{deepseek-v4-flash}, deliberation off    & 6\,s   & 0 \\
    \bottomrule
  \end{tabular}
\end{table}

Twenty-five times faster end to end, for an edit of the same quality. The two
settings that get you there are both on by default:

\begin{itemize}[leftmargin=*]
  \item \textbf{A fast model.} Reach for a heavier one only where the reasoning
        earns its keep, which proofreading a paragraph does not.
  \item \textbf{\ui{Answer directly, without deliberating}}, in
        \ui{Settings} $\rightarrow$ \ui{Speed}. Turn it off when you want the
        model to think a problem through.
\end{itemize}

If a model does deliberate, the panel streams its chain of thought under
\ui{Thinking} so you can see it working rather than hung. That text is never
part of the answer and is never written into your document.

\paragraph{A note on Max Tokens.}
Deliberation is charged against the same output budget as the answer. A model
that thinks for twelve thousand tokens on a four thousand token budget produces
nothing at all --- so the default is 16{,}384, and the editor tells you plainly
when a run ended that way rather than showing an empty result.

\subsection{The inline assistant}
\label{sec:inline}

Most edits should never involve the sidebar at all. Put the cursor in a
paragraph and press \kbd{Cmd+K}.

A bar opens under the text with the passage it is about to change tinted behind
it. You did not have to select anything: with no selection, \kbd{Cmd+K} takes
the paragraph you are standing in, bounded by blank lines and by
\verb|\begin|/\verb|\end|, so pressing it inside a table body works on the row
rather than swallowing the environment. Select first if you want a different
scope.

Then either type what you want --- \emph{say this in plainer language},
\emph{cut the hedging} --- or click one of the presets. The answer streams in,
you get a diff, and:

\begin{itemize}[leftmargin=*]
  \item \kbd{Cmd+Enter} accepts it.
  \item \kbd{Esc} discards it, or stops a run in progress. It works wherever
        the keyboard focus happens to be, including after you have clicked back
        into the document --- and there is a close button on the bar for the
        same purpose.
  \item \kbd{Enter} refines: type \emph{shorter}, \emph{less formal}, and it
        runs again --- against the \emph{original} text, not its own last
        answer, so refinements do not compound edits on top of edits.
\end{itemize}

Nothing is written until you accept, and \kbd{Cmd+Z} undoes an accept like any
other edit.

\subsection{Reading a diff}

Where a line was rewritten rather than replaced wholesale, only the words that
actually changed are tinted --- the rest of the line is dimmed:

\begin{lstlisting}
- are all coloured differently, [\verb|\begin|/\verb|\end|] blocks fold, and
+ are all coloured differently, [code] blocks fold, and
\end{lstlisting}

This matters more than it sounds. A plain line diff shows a one-word fix as a
whole line deleted and a whole line added, and leaves you to find the
difference yourself. Pairing only happens where it helps: if the two versions
share no wording, or the changes are scattered across the whole line, you get
the two plain lines instead, because a word diff there is confetti.

\subsection{The loop}

\begin{enumerate}[leftmargin=*]
  \item \textbf{Select} the passage you want changed.
  \item Open \ui{AI} in the sidebar and pick a skill.
  \item Press \ui{Run}. The answer streams in; the button becomes
        \ui{Cancel} with a running timer.
  \item \textbf{Read the diff.} Removed lines in red, added in green,
        unchanged stretches collapsed to \texttt{$\cdots$ n unchanged lines}.
  \item Press \ui{Apply} --- or do not.
\end{enumerate}

\subsection{Selection is the point}

Before you press \ui{Run}, the panel states what it is about to send:
\emph{selection (4 lines)}, or \emph{whole file (312 lines)}. That line is the
one to check.

With a paragraph selected, only that paragraph goes to the model, and
\ui{Apply} writes the result back over exactly those lines --- the rest of your
document is untouched, and you pay for a paragraph instead of a thesis. With
nothing selected the whole file is sent, which is right for
\ui{Generate Abstract} and wasteful for \ui{Proofread Section}.

The range is captured when you press \ui{Run}, not when you press \ui{Apply}.
You are free to scroll and click around while the model works. And if you
switch to a different file, \ui{Apply} is disabled with a note explaining why
--- rather than writing your revised paragraph into whatever document happens
to be open.

\subsection{Skills do different things}

Skills are not interchangeable, and the panel treats them differently.

\begin{table}[htbp]
  \centering
  \caption{Skill kinds and what the panel offers for each}
  \label{tab:skills}
  \begin{tabular}{lll}
    \toprule
    \textbf{Kind} & \textbf{Skills} & \textbf{You get} \\
    \midrule
    Rewrite & Proofread, Tighten, Improve Equation,   & a diff, and \ui{Apply} \\
            & Fix Errors, Fix Bibliography            & \\[2pt]
    Generate & Generate Table, Paste Data             & the \LaTeX, and \\
             & $\rightarrow$ Table, Generate Abstract, & \ui{Insert at cursor} \\
             & Executive Summary, Notes               & \\
             & $\rightarrow$ Prose, Plain Text
               $\rightarrow$ \LaTeX                   & \\[2pt]
    Explain & Explain Equation, Summarize Section,    & prose, and \ui{Copy} \\
            & Consistency Check,                      & only \\
            & Suggest Paper Structure                 & \\
    \bottomrule
  \end{tabular}
\end{table}

There is no \ui{Apply} button on \ui{Explain Equation}, because an explanation
is for you to read, not to paste into your paper. A diff would be meaningless
there too, so none is shown.

\subsection{It reads your project}

Each skill declares what it needs to know, and the panel gathers it before
sending --- listing what it collected under \emph{Also sends}, so the cost is
never hidden from you.

\begin{table}[htbp]
  \centering
  \caption{Context attached to each skill}
  \label{tab:context}
  \begin{tabular}{ll}
    \toprule
    \textbf{Skill} & \textbf{Also sends} \\
    \midrule
    Proofread Section, Tighten     & labels, citation keys \\
    Inline edit (\kbd{Cmd+K})      & preamble, labels, citation keys \\
    Improve Equation               & preamble \\
    Generate Table, Paste Data
      $\rightarrow$ Table          & preamble \\
    Notes $\rightarrow$ Prose,
      Plain Text $\rightarrow$ \LaTeX & preamble, labels, citation keys \\
    Fix Compilation Errors         & the errors, preamble \\
    Fix Bibliography               & citation keys \\
    Consistency Check              & labels \\
    Suggest Paper Structure        & project file tree \\
    Explain, Summarize,
      Executive Summary            & nothing \\
    \bottomrule
  \end{tabular}
\end{table}

This is what stops the two most irritating failure modes. A model that has been
shown your preamble will not hand you a \verb|\toprule| for a document that
never loaded \texttt{booktabs}. A model that has been given the keys in your
\texttt{.bib} will cite \texttt{knuth1984} rather than inventing
\texttt{smith2019}.

\ui{Explain Equation} sends nothing extra, because nothing extra would help ---
paying for a file tree to explain a formula is waste.

\subsection{Worked example: fixing a compilation error}

This is the skill where the difference is starkest, because the editor already
has everything needed.

Suppose you typed:

\begin{lstlisting}
\includegrpahics[width=0.8\textwidth]{plot.pdf}
\end{lstlisting}

Press \ui{Typeset} and \ui{Issues} reports
\emph{Undefined control sequence}. Now:

\begin{enumerate}[leftmargin=*]
  \item Select the offending line.
  \item \ui{AI} $\rightarrow$ \ui{Fix Compilation Errors} $\rightarrow$ \ui{Run}.
\end{enumerate}

You are not asked to paste the log. The panel shows
\emph{Also sends: 1 error, preamble}, having taken the error straight from the
run you just did. The diff comes back:

\begin{lstlisting}
- \includegrpahics[width=0.8\textwidth]{plot.pdf}
+ \includegraphics[width=0.8\textwidth]{plot.pdf}
\end{lstlisting}

\ui{Apply}, then \ui{Typeset}.

\subsection{Worked example: tightening a paragraph}

Select a paragraph you are unhappy with and run \ui{Proofread Section}. A
typical diff:

\begin{lstlisting}
- The experiments that we ran showed good results across all
- three of the datasets that we evaluated on.
+ Our experiments show strong results across all three datasets.
\end{lstlisting}

Because the skill also receives your labels and citation keys, a suggestion to
cite something will name a key you actually have.

Read every diff before applying. The assistant is a fast, well-informed
copy-editor, not an authority on your argument --- and on a rewrite skill it
will happily reword a sentence whose awkwardness was deliberate.

\subsection{Copy Prompt}

\ui{Copy Prompt} puts the exact prompt --- context included --- on your
clipboard, for pasting into another tool. It is the same text \ui{Run} would
send, so what you get elsewhere matches what you would have got here.

\subsection{What it is actually good at}
\label{sec:ai-measured}

Every section of this guide was run through \ui{Proofread Section}. The results
are worth stating plainly, because they describe the tool better than any
feature list.

\paragraph{It does not break your \LaTeX.}
Across all six sections, every \verb|\begin| matched its \verb|\end|, and every
\verb|\label| and \verb|\ref| came back intact. That is the property that makes
\ui{Apply} safe to press.

\paragraph{On polished prose, it finds almost nothing.}
Six sections produced exactly one correction worth taking --- a missing comma
between two independent clauses:

\begin{lstlisting}
- It does not upload your documents anywhere and it does not
+ It does not upload your documents anywhere, and it does not
\end{lstlisting}

Everything else was house style rather than error: serial commas, added in some
lists and not others; and removing the spaces around em-dashes, which this
document uses consistently, so accepting it would have made one section
disagree with the other five. Applied wholesale, those edits would have left
the guide \emph{less} consistent than it started.

\paragraph{Watch for reflowing.}
The largest ``change'' was not a change at all. On one section the model
returned text that was word-for-word identical to the input, rewrapped onto
different lines --- forty-eight changed lines, no difference in the PDF. The
skills now instruct the model to preserve line breaks, which reduced that to
zero, but it is the failure mode to recognise: a diff that is enormous and
entirely uniform is usually reflowing, not editing.

\paragraph{So what is it for?}
Mechanics, on text that is still rough. A first draft, notes pasted in from
elsewhere, a paragraph written at midnight --- there it earns its place
immediately. On prose you have already revised twice, expect it to propose
preferences and expect to decline most of them. That is not a failure of the
model; it is what running a copy-editor over clean copy looks like.

Which is why every rewrite skill ends in a diff and a button. Judge each
changed line on its own: accepting a whole rewrite because most of it is
harmless is how a document loses its voice one paragraph at a time.

\subsection{When it goes wrong}

If a request fails, the error appears in its own red panel and the result area
stays empty --- an error never masquerades as output you might paste into your
document. \ui{Cancel} stops a run in progress and keeps whatever streamed in so
far, in case the useful part arrived early.

\section{Version control, agents and plugins}
\label{sec:integrations}

\subsection{Git}

The \ui{Git} panel covers the day-to-day loop without leaving the editor: it
shows the current branch, how far ahead or behind you are, and your staged,
modified and untracked files. You can stage everything, write a message and
commit, push, pull, and view the diff of the file you are editing. If the
folder is not a repository yet, the panel offers to initialise one.

It is deliberately not a full Git client. Rebases, conflict resolution and
history surgery belong in a terminal or a dedicated tool. What is here is what
you want between paragraphs.

A \LaTeX{}-specific note: commit your sources, not your output. A
\texttt{.gitignore} containing

\begin{lstlisting}
build/
*.aux
*.log
*.out
*.toc
*.synctex.gz
\end{lstlisting}

keeps regenerable files out of your history. Since the editor writes all output
to \texttt{build/}, the first line does most of the work.

\subsection{The MCP server}

The editor runs a small Model Context Protocol server on
\texttt{127.0.0.1:9876}. Its purpose is to let an external coding agent see and
act on the project you have open --- useful when you want something more
open-ended than a single skill, such as restructuring a document across several
files.

It exposes six tools:

\begin{table}[htbp]
  \centering
  \caption{MCP tools}
  \label{tab:mcp}
  \begin{tabular}{ll}
    \toprule
    \textbf{Tool} & \textbf{Purpose} \\
    \midrule
    \texttt{read\_file}             & read a file from the project \\
    \texttt{write\_file}            & write or overwrite a file \\
    \texttt{compile}                & typeset and return errors and warnings \\
    \texttt{get\_errors}            & parse a \texttt{.log} for its errors \\
    \texttt{get\_project\_structure} & return the file tree \\
    \texttt{get\_section\_structure} & extract headings from a file \\
    \bottomrule
  \end{tabular}
\end{table}

The indicator at the top of the \ui{AI} panel reports whether the server
actually bound its port. If it reads \emph{not running}, something else is
using 9876; nothing else in the editor is affected.

\subsubsection*{Connecting an agent}

The server speaks JSON-RPC over HTTP. Any MCP client can point at it; for
Claude Code, add it as an HTTP server:

\begin{lstlisting}
claude mcp add --transport http latex http://127.0.0.1:9876/mcp
\end{lstlisting}

To check it by hand before wiring anything up, ask it what it can do:

\begin{lstlisting}
curl -s http://127.0.0.1:9876/mcp \
  -H 'Content-Type: application/json' \
  -d '{"jsonrpc":"2.0","id":1,"method":"tools/list"}'
\end{lstlisting}

The project the editor currently has open is what the tools operate on --- the
app pushes the open project and active file to the server as you work, so an
agent always sees what you see.

\subsubsection*{Skills or agent?}

The two AI surfaces are for different shapes of work, and it is worth being
deliberate about which you reach for.

\begin{table}[htbp]
  \centering
  \caption{Choosing between the panel and an external agent}
  \label{tab:ai-surfaces}
  \begin{tabular}{p{0.44\textwidth}p{0.44\textwidth}}
    \toprule
    \textbf{AI Skills panel} & \textbf{External agent over MCP} \\
    \midrule
    One passage, one task, one diff you approve.
      & Many files, several steps, its own plan. \\[3pt]
    You see exactly what is sent and exactly what changes.
      & It reads and writes on its own; review with \texttt{git diff}. \\[3pt]
    Right for: proofreading, fixing an error, generating a table.
      & Right for: restructuring a document, renaming a label everywhere,
        splitting a chapter. \\
    \bottomrule
  \end{tabular}
\end{table}

Because \texttt{write\_file} overwrites without asking, commit before letting an
agent loose --- then the review is a diff rather than an act of faith.

Two cautions. The server is bound to localhost and unauthenticated: anything
able to reach that port can read and write the open project, so do not forward
it off the machine. And \texttt{write\_file} overwrites without asking, so
prefer to have an agent work on a project that is committed to Git.

\subsection{Plugins}

Plugins are loaded from \texttt{\textasciitilde/.latex-editor/plugins/}, each a
folder with a \texttt{plugin.json} manifest and an entry point. The
\ui{Plugins} dialog lists what is installed and lets you inspect each one's
manifest and source.

Being straight about the current state: plugins are \emph{listed and
inspectable, but not yet executed}. The manager is a viewer. Two samples ship
in the repository's \texttt{plugins/} folder to show the intended manifest
shape --- treat the system as a preview of an interface rather than a working
extension point.

\subsection{SyncTeX}

The backend can map a source line to a PDF page and back. It is not yet wired
to a click in the preview, so there is no jump-to-source today. The groundwork
is done; the gesture is not.

\section{Reference}
\label{sec:reference}

\subsection{Keyboard shortcuts}

\begin{table}[htbp]
  \centering
  \caption{Shortcuts (use \kbd{Ctrl} in place of \kbd{Cmd} on Windows and Linux)}
  \label{tab:shortcuts}
  \begin{tabular}{ll}
    \toprule
    \textbf{Keys} & \textbf{Action} \\
    \midrule
    \kbd{Cmd+Enter}       & Typeset \\
    \kbd{Cmd+S}           & Save now (autosave handles this anyway) \\
    \kbd{Cmd+Shift+P}     & Command palette (lists all shortcuts) \\
    \midrule
    \kbd{Cmd+K}           & Inline AI edit on the selection or paragraph \\
    \kbd{Cmd+Enter}       & \dots{} accept the proposed edit \\
    \kbd{Esc}             & \dots{} discard it, or stop a run \\
    \kbd{Enter}           & \dots{} refine and run again \\
    \midrule
    \kbd{Cmd+F}           & Find \\
    \kbd{Cmd+Alt+F}       & Find and replace \\
    \kbd{Cmd+G}           & Go to line \\
    \kbd{Cmd+/}           & Toggle comment \\
    \kbd{Shift+Alt+Down}  & Duplicate line \\
    \kbd{Alt+Up} / \kbd{Alt+Down} & Move line \\
    \midrule
    \kbd{Cmd+}\texttt{+} / \kbd{Cmd+}\texttt{-} & Zoom the preview in / out \\
    \kbd{Cmd+0}           & Preview back to 100\% \\
    \bottomrule
  \end{tabular}
\end{table}

The preview zoom keys apply while the pointer is over the preview. Both panes
are on screen at once, so the shortcut has to mean whichever one you are
looking at --- with the pointer over the editor, the same keys resize the
editor text instead.

Anything not bound to a key is in the command palette, including toggling the
sidebar, preview and issues panel.

\subsection{Zooming the preview}

Three ways, all doing the same thing:

\begin{itemize}[leftmargin=*]
  \item The \ui{$-$} and \ui{$+$} buttons, which step through a ladder of
        round numbers rather than drifting onto values like 137\%.
  \item Pinch on a trackpad, or \kbd{Ctrl}+scroll, which zooms continuously and
        keeps the point under the pointer still --- so you can zoom into a
        figure by pointing at it.
  \item The percentage itself, which opens presets plus \ui{Fit width} and
        \ui{Fit page}.
\end{itemize}

The two fit options are modes, not one-off calculations: drag the splitter and
the page re-fits to the new pane width. Any manual zoom leaves the mode, since
otherwise the next resize would undo what you just did.

Pages render at your display's full pixel density, so zooming in shows more
detail rather than magnifying a blurry image.

\subsection{Choosing a compiler}

\begin{description}[style=nextline, leftmargin=0pt, labelindent=0pt, labelwidth=0pt, font=\normalfont\ttfamily]
  \item[pdflatex] The default. Fastest, and correct for most documents.
  \item[xelatex] Choose when you need system fonts via \texttt{fontspec}, or
        substantial non-Latin script.
  \item[lualatex] Choose for \texttt{fontspec} plus embedded Lua, or when a
        package requires it.
  \item[latexmk] Not a compiler but a driver: it re-runs the real compiler
        until references settle. The one to pick if the two-press rule of
        Section~\ref{sec:passes} irritates you, or if you need BibTeX run
        automatically.
\end{description}

\subsection{Troubleshooting}

\begin{description}[style=nextline, leftmargin=0pt, labelindent=0pt, labelwidth=0pt, font=\normalfont\itshape]

  \item[The compiler menu is empty.]
  No \LaTeX{} distribution was found on your \texttt{PATH}. Install TeX Live,
  MacTeX or MiKTeX and restart. Verify with \texttt{which pdflatex} in a
  terminal --- if the shell cannot find it either, the problem is the
  installation, not the editor.

  \item[References show as \texttt{??}, the contents page is empty.]
  Press \ui{Typeset} a second time. See Section~\ref{sec:passes}.

  \item[Citations are undefined.]
  The editor does not run BibTeX or Biber. Either use a
  \texttt{thebibliography} environment as this guide does, run
  \texttt{biber}/\texttt{bibtex} once in a terminal, or switch to
  \texttt{latexmk}.

  \item[The wrong file is being typeset.]
  The editor builds the file containing \verb|\documentclass|. If your project
  has several, the shortest path wins. Keep one root document per project
  folder.

  \item[The AI panel says ``whole file'' when you wanted a paragraph.]
  Nothing is selected. Select the text first --- the panel always tells you
  what it is about to send.

  \item[\ui{Apply} is greyed out.]
  Either the model proposed no change, or you switched files after running.
  The panel says which.

  \item[The MCP indicator says ``not running''.]
  Another process holds port 9876. Everything else works.

\end{description}

\subsection{Current limitations}

Stated plainly, so nothing here is a surprise:

\begin{itemize}[leftmargin=*]
  \item One file open at a time --- there are no editor tabs.
  \item One compiler pass per press, with the consequences in
        Section~\ref{sec:passes}.
  \item SyncTeX exists in the backend but is not bound to a click.
  \item Plugins are listed, not executed.
  \item AI context collects \verb|\label|s from the open file only; citation
        keys are gathered across the whole project.
  \item One AI request at a time.
\end{itemize}

\subsection{Further reading}

For \LaTeX{} itself rather than this editor, Lamport~\cite{lamport1994} remains
the clearest introduction, Mittelbach and Goossens~\cite{mittelbach2004} is the
reference to own when a package fights you, and Knuth~\cite{knuth1984} is where
all of it comes from.


% The editor runs one compiler pass, so a \bibliography{refs} call
% would never resolve — BibTeX is a separate program the app does not
% invoke. An inline bibliography resolves in the same single pass.
% refs.bib ships alongside anyway: the AI assistant reads it to learn
% your real citation keys (see Section 4).
\begin{thebibliography}{9}

\bibitem{knuth1984}
Donald E. Knuth.
\emph{The \TeX book}.
Addison--Wesley, 1984.

\bibitem{lamport1994}
Leslie Lamport.
\emph{\LaTeX: A Document Preparation System}.
2nd edition. Addison--Wesley, 1994.

\bibitem{mittelbach2004}
Frank Mittelbach and Michel Goossens.
\emph{The \LaTeX\ Companion}.
2nd edition. Addison--Wesley, 2004.

\end{thebibliography}

\end{document}
